% Les primitives constituent les briques de calcul de la bibliothèque.
% Les contributions personnelles sont précisées dans content/03-mission.tex.

\subsection{Agrégateurs}

% MIN 8:00 à 9:00. Présenter les mécanismes et la moyenne de référence.
\begin{frame}{Les familles d'agrégateurs}
    J'ai pris en charge leur implémentation et leurs tests.

    \medskip
    \begin{tabular}{@{}p{0.24\textwidth}p{0.72\textwidth}@{}}
        \textbf{Distances} & Sélection de gradients proches. \\[0.35em]
        \textbf{Coordonnées} & Filtrage de chaque coordonnée séparément. \\[0.35em]
        \textbf{Hybrides} & Combinaison de sélection et de filtrage. \\[0.35em]
        \textbf{Géométrie} & Recherche d'un centre des vecteurs.
    \end{tabular}

    \begin{block}{La moyenne comme référence : \average{}}
        Pour les gradients reçus $\vect{g}_1,\ldots,\vect{g}_{\nworkers}$,
        \[
            \operatorname{Average}=\frac{1}{\nworkers}
            \sum_{i=1}^{\nworkers}\vect{g}_i.
        \]
        La moyenne inclut toutes les contributions, même arbitraires.
    \end{block}
    % À l'oral, rappeler l'appui sur les règles publiées et le code de Sébastien.
\end{frame}

% MIN 9:00 à 10:35. Lire le score, puis la moyenne de la variante implémentée.
% L'annexe du rapport précise l'écart avec la construction itérative publiée.
\begin{frame}{\mkrum{} : sélectionner puis moyenner}
    Au plus $\nbyz$ des $\nworkers$ participants sont \textbf{malhonnêtes}.

    \medskip
    $\mathcal N_i$ contient les $\nworkers-\nbyz-2$ plus proches voisins
    de $\vect{g}_i$ parmi les \textbf{autres} gradients.
    \[
        s_i=\sum_{j\in\mathcal N_i}
        \norm{\vect{g}_i-\vect{g}_j}^{2}.
    \]
    On calcule les scores une seule fois. $I_{\nselected}$ contient les indices
    des $\nselected$ gradients ayant les scores les plus faibles.
    \[
        \operatorname{MultiKrum}(\vect{g}_1,\ldots,\vect{g}_{\nworkers})
        =\frac{1}{\nselected}\sum_{i\in I_{\nselected}}\vect{g}_i.
    \]
    Pour $\nselected=1$, on retrouve \textbf{\krum{}}.
    Pour $\nselected=n$, on retrouve \textbf{Average}.
\end{frame}

\subsection{Attaques}

% MIN 10:35 à 11:25. Expliquer les deux attaques utilisées dans les résultats.
\begin{frame}{Deux attaques pour éprouver les agrégateurs}
    \alert{\textbf{Rappel : scénario d'omniscience.}}\\
    Les attaquants ont accès aux gradients honnêtes.

    \medskip
    On note $\bar{\vect{g}}$ la moyenne des gradients honnêtes.
    \begin{columns}[T,onlytextwidth]
        \begin{column}{0.47\textwidth}
            \begin{block}{\signflip{}}
                L'attaque inverse et amplifie la moyenne des gradients honnêtes.
                \[
                    \vect{b}=-\lambda\bar{\vect{g}},\qquad\lambda>0.
                \]
                Tous les participants malhonnêtes envoient ce même vecteur.
            \end{block}
        \end{column}
        \begin{column}{0.49\textwidth}
            \begin{block}{\smallperturbation{}}
                L'attaque décale une seule coordonnée de la moyenne honnête.
                \[
                    \vect{b}=\bar{\vect{g}}+\gamma\vect{e}_k.
                \]
                $\vect{e}_k$ est le vecteur de base de la coordonnée $k$.
                Le vecteur cherche à rester sélectionnable.
            \end{block}
        \end{column}
    \end{columns}
    % La coordonnée de SmallPerturbation a la plus forte variance honnête.
    % L'expérience de 2018 utilise une amplitude fixée, présentée avec ses réglages.
\end{frame}

\subsection{Partitionnement des données}

% MIN 11:25 à 12:10. Distinguer hétérogénéité honnête et attaque.
\begin{frame}{Les données locales peuvent aussi différer}
    \begin{columns}[T,onlytextwidth]
        \begin{column}{0.47\textwidth}
            \begin{block}{Données IID}
                Observations indépendantes et de même loi.

                \smallskip
                Répartition aléatoire en sous-jeux comparables.
            \end{block}
        \end{column}
        \begin{column}{0.49\textwidth}
            \begin{block}{Données non IID}
                Les lois locales peuvent différer.

                \smallskip
                Exemple : les hôpitaux accueillent des populations différentes.
            \end{block}
        \end{column}
    \end{columns}
    \begin{block}{Partitionneurs}
        Distribuer les données entre participants et régler leur hétérogénéité
        pour comparer les agrégateurs.
    \end{block}
    \alert{Même sans attaque, les gradients honnêtes peuvent différer.}
\end{frame}

\subsection{Extension des primitives}

% MIN 12:10 à 12:55. Montrer comment ajouter une primitive réutilisable.
\begin{frame}{Une seule méthode pour ajouter une règle}
    Une méthode est une fonction associée à un composant logiciel.

    \medskip
    \begin{tabular}{@{}ll@{}}
        \textbf{Primitive} & \textbf{Méthode à implémenter} \\[0.4em]
        Agrégateur & \code{aggregate} \\[0.3em]
        Attaque & \code{generate} \\[0.3em]
        Partitionneur & \code{partition}
    \end{tabular}

    L'interface commune permet de réutiliser la primitive dans les simulations
    sans réécrire leur boucle.

    \medskip
    Au total, la bibliothèque comprend onze règles, cinq attaques,
    quatre partitionneurs et six architectures de modèles.
\end{frame}

\subsection{Modèles pour l'agrégation}

% MIN 12:55 à 15:00. Expliquer le rôle du composant sans détailler l'implémentation.
% Le dessin utilise picture et les couleurs Beamer, sans dépendance supplémentaire.
\begin{frame}{Préparer les modèles pour l'agrégation}
    PyTorch stocke les paramètres et les gradients par couche.
    Pour les agréger, le composant \code{Model} les présente sous forme
    d'une suite de valeurs.

    \begin{center}
        \begingroup
        \setlength{\unitlength}{1mm}
        \setlength{\fboxsep}{0pt}
        \begin{picture}(112,33)
            \put(0,25){\makebox(39,6)[l]{Valeurs par couche}}
            \put(42,23){\colorbox{structure.fg!15}{\makebox(44,10){$\begin{matrix}a&b\\c&d\end{matrix}$}}}
            \put(86,23){\colorbox{defenseAccent!20}{\makebox(22,10){$[\,e\quad f\,]$}}}
            \put(64,22){\vector(0,-1){5}}
            \put(64,17){\vector(0,1){5}}
            \put(97,22){\vector(0,-1){5}}
            \put(97,17){\vector(0,1){5}}
            \put(0,10){\makebox(39,6)[l]{Valeurs réunies}}
            \put(42,10){\colorbox{structure.fg!15}{\makebox(11,6){$a$}}}
            \put(53,10){\colorbox{structure.fg!15}{\makebox(11,6){$b$}}}
            \put(64,10){\colorbox{structure.fg!15}{\makebox(11,6){$c$}}}
            \put(75,10){\colorbox{structure.fg!15}{\makebox(11,6){$d$}}}
            \put(86,10){\colorbox{defenseAccent!20}{\makebox(11,6){$e$}}}
            \put(97,10){\colorbox{defenseAccent!20}{\makebox(11,6){$f$}}}
            \put(42,7){\line(1,0){66}}
            \put(42,7){\line(0,1){2}}
            \put(108,7){\line(0,1){2}}
            \put(42,0){\makebox(66,4){Même mémoire}}
        \end{picture}
        \endgroup
    \end{center}
    La vue réunie partage la mémoire avec les valeurs par couche, ce qui évite
    de les recopier pour l'agrégateur. Le calcul des gradients reste nécessaire.
    % Le même principe s'applique aux paramètres et aux gradients.
    % Le schéma représente l'un de ces ensembles, avec son propre bloc de mémoire.
    % Une concaténation naïve recopierait le vecteur à chaque accès.
\end{frame}
